% !TITLE 葬花吟（节选）
% !AUTHOR 曹雪芹
% !DYNASTY 明清
% !GENRE 七言古诗
% !ORDER 6
% !SOURCE 《红楼梦》（借林黛玉之手）
% !EDITION 《红楼梦》，中国艺术研究院红楼梦研究所校注，人民文学出版社（2008年版为候选）
% !EDITION_STATUS pending

\begin{poem}
花谢花飞花满天，红消香断\textsuperscript{1}有谁怜？
游丝\textsuperscript{2}软系飘春榭\textsuperscript{3}，落絮\textsuperscript{4}轻沾扑绣帘。
闺中女儿惜春暮，愁绪满怀无释处。
手把花锄出绣帘，忍踏落花来复去。

……

一年三百六十日，风刀霜剑\textsuperscript{5}严相逼。
明媚鲜妍能几时，一朝漂泊难寻觅。

……

未若锦囊\textsuperscript{6}收艳骨\textsuperscript{7}，一抔净土\textsuperscript{8}掩风流\textsuperscript{9}。
质本洁来还洁去，强于污淖\textsuperscript{10}陷渠沟。

……

侬\textsuperscript{11}今葬花人笑痴，他年葬侬\textsuperscript{12}知是谁？
试看春残花渐落，便是红颜老死时。
一朝春尽红颜老，花落人亡两不知！
\end{poem}

\begin{annotations}
\notetext{1}{红消香断：花的红色褪去、香气断绝。指花瓣凋谢枯萎。}
\notetext{2}{游丝：春天飘在空中的蛛丝或虫丝。蜘蛛吐丝在风中游荡，故称“游丝”。}
\notetext{3}{春榭：春天的台榭（建在高台上的敞屋）。游丝飘过春天的亭台。}
\notetext{4}{落絮：飘落的柳絮。柳絮轻飞，沾在绣花的门帘上。}
\notetext{5}{风刀霜剑：风像刀，霜像剑。形容环境严酷，岁月对人的摧残。}
\notetext{6}{锦囊：锦缎做的小袋子。林黛玉用锦囊收集落花安葬。}
\notetext{7}{艳骨：艳丽的花骨——指花瓣。“艳”写花之美，“骨”写花已死，美与死的并置。}
\notetext{8}{一抔净土：一捧干净的泥土。抔（póu），用手捧。}
\notetext{9}{风流：风韵、美好姿态。埋入土中的不只是花瓣，还有花的风韵。}
\notetext{10}{污淖：污浊的泥沼。淖（nào），烂泥。与其被抛弃在肮脏的泥沟里，不如让它在净土中安眠。}
\notetext{11}{侬：我。吴语方言（苏州话）中“我”的说法。林黛玉是苏州人，故用此自称。}
\notetext{12}{他年葬侬：将来埋葬我。“他年”即将来。林黛玉由葬花想到自己的结局——今天花落了有我埋葬，将来我死了谁来埋葬我呢？}
\end{annotations}

\begin{appreciation}
《葬花吟》是《红楼梦》第二十七回里林黛玉所吟的一首长诗，也是中国文学史上写怜花与自怜最动人的诗篇之一。

先要知道它是在什么情形下吟出来的。芒种那一天，大观园里众女儿忙着祭饯花神，嬉笑热闹；前一夜，黛玉去怡红院叩门，晴雯没听出她的声音把她关在门外，她只当是宝玉恼了她。第二天满园欢笑，她独自扛着花锄走到花冢边，一边葬花，一边落泪。越是热闹的日子她越显得孤零零的。懂了这一层，才读得懂这首诗。

花谢花飞花满天，一句里三个花字，满眼都是落花；红消香断有谁怜，紧接一个问句，答案是没有人。花之美正在消失，而消失无人过问。游丝软系飘春榭，落絮轻沾扑绣帘——蛛丝软软地系在亭台边柳絮轻轻扑上帘子，两句都写一个轻字，写的都是留不住的柔物。手把花锄出绣帘，忍踏落花来复去，一个忍字最细：不忍心踩下去却又不能不挪动脚步，只好绕着落花来回躲着走。

一年三百六十日，风刀霜剑严相逼——春风春雨，在旁人眼里是滋养在她眼里是刀是剑。明媚鲜妍能几时，一朝漂泊难寻觅，花如此，人也如此。这就逼出了全诗的骨：未若锦囊收艳骨，一抔净土掩风流。质本洁来还洁去，强于污淖陷渠沟。与其让花瓣在渠沟里烂掉，不如收进锦囊、埋进净土；干干净净地来，干干净净地去——这几乎是黛玉替自己的一生先写下的一句判词。

侬今葬花人笑痴，他年葬侬知是谁？侬是苏州话里的我。黛玉是苏州人，最伤心的时候，她用回了家乡的口语。今天花落了有我埋葬将来我死了谁来埋葬我？《氓》里那个被弃的女子也是这样孤单：兄弟不知，咥其笑矣。古老的诗里她被人笑后来的小说里也一样。到最后一朝春尽红颜老，花落人亡两不知，花与人一并消失，再也无从相认。

怜花的题目早有人做过。刘希夷写年年岁岁花相似岁岁年年人不同，是把花期与人生对着看；《桃夭》里桃之夭夭灼灼其华是花开的祝福。黛玉却把花的结局直接搬到自己身上——她葬的是花，想到的始终是自己。前面读过唐寅的《桃花庵歌》，唐寅也爱花，他的花是陪他喝酒的伴儿，落就落了，无花无酒锄作田，落得潇洒；黛玉的花是她自己的化身，落一朵心尖上便疼一下。同一个春天，两种活法。

黛玉此前在梨香院外偶然听见女伶唱《牡丹亭》的良辰美景奈何天，一时心动神摇，几乎站不住。那时她听见的是别人的词；到芒种这一天，她把那一句触动，写成了自己的一整首诗。《红楼梦》最了不起的地方之一，就是让一个十几岁的女孩子，把自己对生死的预感一字一字写成了诗。
\end{appreciation}
